% ==========================================
% Week 7: DIFFUSION PDEs
% ==========================================
\section{Week 7 - Diffusion and Partial Differential Equations}

\subsection{Non-Homogeneous Boundary Conditions (BC)}

\textbf{Concept:} The idea is to decompose the full solution into a steady-state part that satisfies the non-homogeneous boundary conditions and a transient part that satisfies homogeneous boundary conditions. As $t \to \infty$, the transient contribution decays and the solution converges to the steady-state profile.

The governing one-dimensional diffusion (or heat) equation is:
\begin{equation}
    \frac{\partial C}{\partial t} = D \frac{\partial^2 C}{\partial x^2}
\end{equation}

\textbf{Boundary Conditions (BC):}
\begin{itemize}
    \item $C(0,t)=C_1$
    \item $C(L,t)=C_2$
\end{itemize}

\textbf{Initial Condition (IC):}
\begin{itemize}
    \item $C(x,0)=C_0(x)$
\end{itemize}

\subsubsection{Step A: Solve for the Steady State}

Define the steady-state solution by
\begin{equation}
    C_f(x)=\lim_{t\to\infty} C(x,t)
\end{equation}

At steady state, the time derivative vanishes:
\begin{equation}
    \frac{\partial C_f}{\partial t}=0
    \qquad \Longrightarrow \qquad
    0 = D\frac{\partial^2 C_f}{\partial x^2}
\end{equation}
Hence,
\begin{equation}
    C_f(x)=C_1+\frac{C_2-C_1}{L}x
\end{equation}

\begin{mainbox}{Critical Analysis: Limits of the Steady-State Approximation}
\textbf{Mathematical rigor vs. physical reality:} The formal definition $C_f(x)=\lim_{t\to\infty}C(x,t)$ refers to an infinite-time limit. In practice, a system is considered close to steady state when the Fourier number $Fo=\frac{Dt}{L^2}$ becomes sufficiently large, typically $Fo>0.2$ as a rough engineering estimate. In addition, the linear profile above assumes that the diffusivity $D$ is constant. If $D$ depends on concentration, temperature, or composition, then the PDE becomes non-linear and the steady-state profile is no longer necessarily linear.
\end{mainbox}

\subsubsection{Step B: Define the Transient Component}

Introduce the deviation from steady state:
\begin{equation}
    \tilde{C}(x,t):=C(x,t)-C_f(x)
\end{equation}

Then $\tilde{C}$ satisfies the homogeneous PDE
\begin{equation}
    \frac{\partial \tilde{C}}{\partial t}=D\frac{\partial^2 \tilde{C}}{\partial x^2}
\end{equation}

The transformed homogeneous boundary conditions are:
\begin{itemize}
    \item $\tilde{C}(0,t)=C(0,t)-C_f(0)=C_1-C_1=0$
    \item $\tilde{C}(L,t)=C(L,t)-C_f(L)=C_2-C_2=0$
\end{itemize}

The transformed initial condition is:
\begin{itemize}
    \item $\tilde{C}(x,0)=C_0(x)-C_f(x)$
\end{itemize}

\subsubsection{Step C: Superposition}

The full solution is reconstructed as
\begin{equation}
    C(x,t)=C_f(x)+\tilde{C}(x,t)
\end{equation}

\subsection{Time-Dependent Diffusion (1D Series Expansion)}

The governing PDE for one-dimensional time-dependent diffusion is:

\begin{defbox}{The Diffusion Equation}
\begin{equation}
    \frac{\partial C}{\partial t}=D\frac{\partial^2 C}{\partial x^2}
\end{equation}
\end{defbox}

\noindent \textbf{Boundary Conditions (BC): Fixed concentration} \\
Assume zero concentration at both boundaries:
\begin{equation}
    C(0,t)=C(L,t)=0
\end{equation}

\begin{notebox}[Superposition Principle]
The PDE is linear and homogeneous, and the boundary conditions are also homogeneous. Therefore, if $C_1(x,t)$ and $C_2(x,t)$ are solutions, then
\begin{equation*}
    C_1(x,t)+C_2(x,t)
    \qquad \text{and} \qquad
    \beta C_1(x,t), \quad \beta\in\mathbb{R}
\end{equation*}
are also solutions.
\end{notebox}

\noindent \textbf{Why is this true? (Superposition Principle)}
\begin{align*}
    \text{PDE: } \quad
    \frac{\partial}{\partial t}(C_1+C_2)-D\frac{\partial^2}{\partial x^2}(C_1+C_2)
    &=
    \left(\frac{\partial C_1}{\partial t}-D\frac{\partial^2 C_1}{\partial x^2}\right)
    +
    \left(\frac{\partial C_2}{\partial t}-D\frac{\partial^2 C_2}{\partial x^2}\right) \\
    &=0+0=0
\end{align*}
and the boundary conditions are preserved because
\begin{align*}
    (C_1+C_2)(0,t) &= C_1(0,t)+C_2(0,t)=0+0=0 \\
    (C_1+C_2)(L,t) &= C_1(L,t)+C_2(L,t)=0+0=0
\end{align*}

\subsubsection{Solution Approach (Linearity)}

\textbf{Step 1: General solution representation} \\
Because of linearity, the solution can be written as a superposition of basis functions:
\begin{equation}
    C(x,t)=\sum_{n=1}^{\infty} A_n \varphi_n(x,t)
\end{equation}
where the coefficients $A_n$ are determined by the initial condition.

\vspace{0.5cm}
\textbf{Step 2: Separation of variables} \\
Assume
\begin{equation}
    C(x,t)=X(x)T(t)
\end{equation}
Substituting into the PDE gives
\begin{equation*}
    X(x)\frac{dT(t)}{dt}=D\,T(t)\frac{d^2X(x)}{dx^2}
\end{equation*}
Dividing by $D\,X(x)T(t)$ yields
\begin{equation}
    \underbrace{\frac{1}{D}\frac{T'(t)}{T(t)}}_{\text{depends only on } t}
    =
    \underbrace{\frac{X''(x)}{X(x)}}_{\text{depends only on } x}
    =-\lambda
\end{equation}
Since the left-hand side depends only on $t$ and the right-hand side only on $x$, both must be equal to a constant $-\lambda$.

This leads to two coupled ODEs:
\begin{align}
    T'(t)+D\lambda T(t) &= 0 \\
    X''(x)+\lambda X(x) &= 0
\end{align}

\subsubsection{Solving the Spatial Problem}

We now solve
\begin{equation}
    X''(x)+\lambda X(x)=0
\end{equation}
subject to
\begin{equation}
    X(0)=X(L)=0
\end{equation}

\begin{defbox}{Evaluating the different cases for $\lambda$}
\textbf{Case I: $\lambda<0$} \\
The solution has the form
\begin{equation*}
    X(x)=\alpha_1 e^{\sqrt{-\lambda}\,x}+\alpha_2 e^{-\sqrt{-\lambda}\,x}
\end{equation*}
Applying the boundary conditions forces $\alpha_1=\alpha_2=0$, giving only the trivial solution.

\vspace{0.2cm}
\textbf{Case II: $\lambda=0$} \\
The solution becomes
\begin{equation*}
    X(x)=\alpha_1 x+\alpha_2
\end{equation*}
Applying the boundary conditions again yields only the trivial solution.

\vspace{0.2cm}
\textbf{Case III: $\lambda>0$} \\
The solution is
\begin{equation*}
    X(x)=\alpha_1 \cos(\sqrt{\lambda}\,x)+\alpha_2 \sin(\sqrt{\lambda}\,x)
\end{equation*}
Applying the boundary conditions gives
\begin{align*}
    X(0) &= \alpha_1 = 0 \\
    X(L) &= \alpha_2 \sin(\sqrt{\lambda}\,L)=0
\end{align*}
To avoid the trivial solution, we require
\begin{equation*}
    \sin(\sqrt{\lambda}\,L)=0
\end{equation*}
which implies
\begin{equation*}
    \sqrt{\lambda}\,L = k\pi,
    \qquad k\in\mathbb{Z}^{+}
\end{equation*}
\end{defbox}

Therefore, the eigenvalues are
\begin{equation}
    \lambda_k=\left(\frac{k\pi}{L}\right)^2,
    \qquad k\in\{1,2,3,\dots\}
\end{equation}
and the corresponding eigenfunctions are
\begin{equation}
    X_k(x)=\sin\left(\frac{k\pi}{L}x\right)
\end{equation}

\textit{Note:} The amplitude factor can be absorbed into the global coefficient $A_k$.

\subsubsection{Solving the Temporal Problem}

For a fixed eigenvalue
\begin{equation}
    \lambda_n=\left(\frac{n\pi}{L}\right)^2
\end{equation}
the time equation becomes
\begin{equation*}
    T_n'(t) = -D\lambda_n T_n(t)
\end{equation*}
whose solution is
\begin{equation}
    T_n(t)=\beta_n e^{-D\left(\frac{n\pi}{L}\right)^2 t}
\end{equation}

\subsubsection{Combining the Spatial and Temporal Solutions}

The basis solutions are therefore
\begin{equation}
    \varphi_n(x,t)=e^{-D\left(\frac{n\pi}{L}\right)^2 t}\sin\left(\frac{n\pi}{L}x\right)
\end{equation}
and the general solution satisfying the homogeneous boundary conditions is
\begin{equation}
    C(x,t)=\sum_{n=1}^{\infty} A_n e^{-D\left(\frac{n\pi}{L}\right)^2 t}\sin\left(\frac{n\pi}{L}x\right)
\end{equation}

\subsubsection{Step 3: Determine the coefficients $\{A_n\}$ from the initial condition}

At $t=0$,
\begin{equation}
    C(x,0)=C_0(x)=\sum_{n=1}^{\infty} A_n \sin\left(\frac{n\pi}{L}x\right)
\end{equation}

To determine the coefficients, multiply both sides by $\sin\left(\frac{m\pi}{L}x\right)$ and integrate from $0$ to $L$:
\begin{equation*}
    \int_0^L C_0(x)\sin\left(\frac{m\pi}{L}x\right)\,dx
    =
    \sum_{n=1}^{\infty}
    A_n
    \int_0^L
    \sin\left(\frac{n\pi}{L}x\right)
    \sin\left(\frac{m\pi}{L}x\right)\,dx
\end{equation*}

\begin{notebox}[Orthogonality Property]
\begin{equation*}
    \int_0^L
    \sin\left(\frac{n\pi}{L}x\right)
    \sin\left(\frac{m\pi}{L}x\right)\,dx
    =
    \begin{cases}
        0 & \text{if } n\neq m \\
        \dfrac{L}{2} & \text{if } n=m
    \end{cases}
\end{equation*}

\textit{Trigonometric identity used:}
\begin{equation*}
    \sin(a)\sin(b)=\frac{1}{2}\big[\cos(a-b)-\cos(a+b)\big]
\end{equation*}
\end{notebox}

By orthogonality, only the term $n=m$ remains:
\begin{equation*}
    \int_0^L C_0(x)\sin\left(\frac{n\pi}{L}x\right)\,dx
    =
    A_n \frac{L}{2}
\end{equation*}
Hence,
\begin{equation}
    A_n=\frac{2}{L}\int_0^L C_0(x)\sin\left(\frac{n\pi}{L}x\right)\,dx
\end{equation}

\vspace{0.5cm}
\begin{resultboxnamed}{Final Full Solution}
For any initial condition $C_0(x)$ that admits a Fourier sine expansion, the complete solution is
\begin{equation}
    C(x,t)=
    \sum_{n=1}^{\infty}
    \left[
        \frac{2}{L}
        \int_0^L
        C_0(\tilde{x})
        \sin\left(\frac{n\pi}{L}\tilde{x}\right)\,d\tilde{x}
    \right]
    e^{-D\left(\frac{n\pi}{L}\right)^2 t}
    \sin\left(\frac{n\pi}{L}x\right)
\end{equation}
\end{resultboxnamed}
\subsection{Visualizing Flow Patterns}

Fluid kinematics describes motion without yet invoking conservation laws. The three main geometric ways to represent a flow are streamlines, pathlines, and streaklines.

\begin{enumerate}
    \item \textbf{Streamlines:} Curves tangent to the instantaneous velocity field at a given time.
    \begin{equation}
        \frac{dy}{dx} = \frac{v}{u}, \qquad 
        \frac{dz}{dx} = \frac{w}{u}, \qquad 
        \frac{dz}{dy} = \frac{w}{v}
        \quad \Longrightarrow \quad
        \frac{d\vec{r}_{SL}}{ds} \times \vec{V} = \vec{0}
    \end{equation}

    \item \textbf{Pathlines:} The actual trajectory followed by one specific fluid particle through time.
    \begin{equation}
        \frac{d\vec{r}_p}{dt} = \vec{V}(\vec{r}_p(t), t)
        \qquad \Longrightarrow \qquad
        \frac{dx}{dt} = u, \ \frac{dy}{dt} = v, \ \frac{dz}{dt} = w
    \end{equation}

    \item \textbf{Streaklines:} The locus of all particles that have passed through a given point $\vec{r}_0$ at previous times $\tau < t$.
\end{enumerate}

\textit{Key Result:} In steady flow, where $\partial \vec{V}/\partial t = 0$, streamlines, pathlines, and streaklines coincide.

\subsubsection{Example: The Sprinkler (Unsteady Case)}

Consider the unsteady velocity field
\begin{equation}
    \vec{V} = U_0 \sin\!\left[\omega\!\left(t - \frac{y}{v_0}\right)\right]\hat{i} + v_0 \hat{j}
\end{equation}

\begin{itemize}
    \item \textbf{Streamlines at $t=0$:}  
    \begin{equation}
        x(y) = \frac{U_0}{\omega}\left(\cos\!\left(\frac{\omega y}{v_0}\right) - 1\right)
    \end{equation}
    so the streamline shape is sinusoidal.

    \item \textbf{Pathline for a particle starting at the origin at $t=0$:}  
    \begin{equation}
        x(t)=0, \qquad y(t)=v_0 t
    \end{equation}
    which is a straight vertical line.
\end{itemize}

This example shows that in unsteady flow, streamlines and pathlines are generally different objects.

\subsection{Acceleration Field}

The acceleration of a fluid particle is the total derivative of the velocity field:
\begin{equation}
    \vec{a} = \frac{D\vec{V}}{Dt}
    = \underbrace{\frac{\partial \vec{V}}{\partial t}}_{\text{Local acceleration}}
    + \underbrace{(\vec{V}\cdot\vec{\nabla})\vec{V}}_{\text{Convective acceleration}}
\end{equation}

\begin{mainbox}{Technical Analysis: The Material Derivative}
The operator $\frac{D}{Dt}$ is fundamental in fluid mechanics because it gives the rate of change experienced by a moving fluid particle. The convective term $(\vec{V}\cdot\vec{\nabla})\vec{V}$ is non-linear and is a major source of complex flow behavior, including mechanisms that contribute to turbulence. It measures the change in velocity caused by a particle moving into a region where the flow field is different, even when the flow is steady in time.
\end{mainbox}

\subsubsection{Example: The Diffuser}

For a steady one-dimensional flow
\begin{equation}
    u(x)=v_0\left(1-\frac{x}{L}\right)
\end{equation}
the local acceleration term vanishes because the flow is steady, so only the convective contribution remains:
\begin{equation}
    a_x = u\frac{\partial u}{\partial x}
    = v_0\left(1-\frac{x}{L}\right)\left(-\frac{v_0}{L}\right)
    = -\frac{v_0^2}{L}\left(1-\frac{x}{L}\right)
\end{equation}

The negative sign indicates deceleration, which is consistent with the widening geometry of a diffuser.

\subsection{Velocity Fields: Eulerian vs. Lagrangian Descriptions}

Once the motion has been visualized geometrically, it is useful to introduce the two standard ways of describing it mathematically.

\textbf{Eulerian Description:} Focuses on fixed points in space and describes the velocity as a field:
\begin{equation}
    \vec{V}(\vec{r}, t) = u\hat{i} + v\hat{j} + w\hat{k}
\end{equation}

\textbf{Lagrangian Description:} Follows individual fluid particles and describes their motion from an initial position:
\begin{equation}
    \vec{V}(\vec{r}_0, t)
    \qquad \text{where } \vec{r}_0 \text{ is the particle's initial position}
\end{equation}

\begin{mainbox}{Structural Foundation: Eulerian vs. Lagrangian Comparison}
Eulerian formulations (Leonhard Euler, 1757) are the standard framework for most continuum fluid mechanics and Computational Fluid Dynamics (CFD), because they describe the flow as fields defined everywhere in space. Lagrangian formulations (Joseph-Louis Lagrange, 1788) are indispensable when one must follow individual particles, droplets, interfaces, or solid bodies.
\vspace{0.2cm}

\begin{tabular}{@{}lll@{}}
\toprule
\textbf{Feature} & \textbf{Eulerian (Field)} & \textbf{Lagrangian (Particle)} \\
\midrule
\textbf{Independent Variables} & Position $\vec{r}$, Time $t$ & Initial Position $\vec{r}_0$, Time $t$ \\
\textbf{Mass Tracking} & Flux through a fixed volume & Constant identity of mass \\
\textbf{Differentiation} & Requires the Material Derivative & Standard time derivative \\
\bottomrule
\end{tabular}
\end{mainbox}